

La práctica tuvo como objetivo determinar experimentalmente la relación entre la constante de Planck y la carga elemental, $h/e$, a partir del efecto fotoeléctrico observado en diodos emisores de luz, LED, de distintas longitudes de onda. Para ello se midió el voltaje ánodo-cátodo, $V_{A-K}$, asociado al encendido de cada LED cuando la corriente que circulaba por él alcanzaba un valor de umbral del orden de algunos microamperes.\\

Se conecto la fuente de voltaje variable de corriente directa en serie con una resistencia de $8.2\,\mathrm{k\Omega}$, un microamperímetro y el LED bajo estudio. El voltímetro se conectó en paralelo con el LED. Se estudiaron 8 LEDs de distinto color: Infrarojo, Rojo, Naranja, Amarillo, Verde, Azul, Violeta y Ultravioleta, los ocho LEDs estuvieron en una sola tablilla dentro de una camara oscura, dentro de la camra oscura, frente al LED estudiado, se coloco tambien un detector CCD lineal Thorlabs LC1-USB con rango espectral de $400$-$1000$nm.\\

Despues de realizar el montaje, se aumento el voltaje de la fuente hasta observar un aumento de la intensidad de la luz detectada por la CCD, en este caso, se registraron simultáneamente la corriente $I_{A-K}$ y el voltaje ánodo-cátodo $V_{A-K}$. Este procedimiento se repitió para cada uno de los ocho LEDs disponibles.\\

Cada LED emitía luz con una longitud de onda característica $\lambda$. A partir de dicha longitud de onda se calculó la frecuencia correspondiente mediante

\begin{equation}
    \nu = \frac{c}{\lambda}
\end{equation}

donde $c$ es la velocidad de la luz en el vacío. Bajo la aproximación de que la energía eléctrica entregada a los portadores de carga se transforma en energía luminosa emitida por el LED, se consideró la relación

\begin{equation}
    eV_{A-K} = h\nu
\end{equation}

De esta expresión se obtiene

\begin{equation}
   \nu = \frac{e}{h} V_{A-K}
\end{equation}

por lo que la pendiente de la gráfica corresponde a $e/h$. Con el conjunto de datos obtenidos se construyó una gráfica de $V_{A-K}$ en función de $\nu$ y se realizó un ajuste lineal de la forma

\begin{equation}
    \nu = m V_{A-K} + b
\end{equation}

donde $m$ es la pendiente de la recta y $b$ la ordenada al origen, de modo que $h/e= 1/m$.\\


Se estudió el espectro de emisión de ocho diodos emisores de luz
(LED): infrarrojo, rojo, naranja, amarillo, verde, azul, violeta y
ultravioleta. Cada LED se encendió individualmente y se colocó frente a
la entrada de un espectrómetro óptico. Para cada fuente se
registró la intensidad luminosa en función de la longitud de onda en nano-metros.\\

La longitud de onda del LED, denotada por $\lambda_{\max}$, se tomó como
la longitud de onda correspondiente al dato experimental con la mayor
intensidad. Para estimar el ancho espectral se empleó el ancho completo a media
altura. Considerando que la intensidad de fondo podía aproximarse
como cero, el nivel de media altura se definió como

\begin{equation}
    I_{1/2}=\frac{I_{\max}}{2}.
\end{equation}

Se determinaron las longitudes de onda $\lambda_L$ y $\lambda_R$ para
las cuales la intensidad alcanzó la mitad del máximo, a la izquierda y
a la derecha del pico, respectivamente. Debido a que el valor
$I_{\max}/2$ no necesariamente coincidió exactamente con uno de los
puntos experimentales, las posiciones de los cruces se calcularon
mediante interpolación lineal entre los dos puntos consecutivos que
encerraban dicho nivel. Para dos puntos $(\lambda_i,I_i)$ y
$(\lambda_{i+1},I_{i+1})$, la longitud de onda correspondiente a la
media altura se calculó mediante

\begin{equation}
    \lambda_{1/2}
    =
    \lambda_i+
    \frac{I_{\max}/2-I_i}
         {I_{i+1}-I_i}
    \left(\lambda_{i+1}-\lambda_i\right).
\end{equation}

Una vez obtenidos ambos cruces, el ancho completo a media altura se
calculó como

\begin{equation}
    \mathrm{FWHM}=\lambda_R-\lambda_L.
\end{equation}

De acuerdo con el criterio adoptado para este experimento, la
incertidumbre asociada con la longitud de onda del LED se tomó como el
semi-ancho a media altura:

\begin{equation}
    \Delta\lambda
    =
    \frac{\mathrm{FWHM}}{2}
    =
    \frac{\lambda_R-\lambda_L}{2}.
\end{equation}

Finalmente, las longitudes de onda se reportaron en la forma

\begin{equation}
    \lambda_{\mathrm{LED}}
    =
    \lambda_{\max}\pm\Delta\lambda.
\end{equation}

Debe señalarse que esta cantidad representa principalmente la anchura
espectral de la emisión del LED y se utilizó como una estimación
operacional de la incertidumbre de la longitud de onda, de acuerdo con
el criterio establecido para la práctica.

